\icode{N16Q-EMI-SI-L3-261004-001}
\nceeID{2016海南}{unkown}
如图，一圆形金属环与两固定的平行长直导线在同一竖直平面内，环的圆心与两导线的距离相等，环的直径小于两导线的间距。两导线中通有大小相等、方向向下的恒定电流，以下说法正确的是\paren[]

\begin{tikzpicture}[decoration={markings,mark=at position 0.5 with {\arrow{>}}},font=\scriptsize]
    \draw (0,0) circle (.5);
    \draw[postaction={decorate},thick] (-1.2,1.2)--++(0,-2.4);
    \draw[postaction={decorate},thick] (1.2,1.2)--++(0,-2.4);
\end{tikzpicture}
\begin{choices}
	\item 金属环向上运动时，则环上的感应电流方向为顺时针方向
	\item 金属环向下运动时，则环上的感应电流方向为顺时针方向
	\item 金属环向左侧直导线靠近时，则环上的感应电流方向为逆时针方向
	\item 金属环向右侧直导线靠近时，则环上的感应电流方向为逆时针方向
\end{choices}
\begin{answer}
    D\qquad 提示：

    \begin{tikzpicture}[decoration={markings,mark=at position 0.5 with {\arrow{>}}},font=\scriptsize]
        \draw[dashed,cyan] (0,1.2)--++(0,-2.4);
        \foreach \x in {-.8,-.4}
            \foreach \y in {-1,-.6,...,1}
                \node at(\x,\y)[teal]{\( \cdot \)};
        \foreach \x in {.8,.4}
            \foreach \y in {-1,-.6,...,1}
                \node at(\x,\y)[teal]{\(\times\)};
	    \draw (0,0) circle (.5);
	    \draw[postaction={decorate},thick] (-1.2,1.2)--++(0,-2.4);
	    \draw[postaction={decorate},thick] (1.2,1.2)--++(0,-2.4);
    \end{tikzpicture}
\end{answer}